\documentclass[11pt,a4paper]{article}
\usepackage[T1]{fontenc}
\usepackage[utf8]{inputenc}
\usepackage{lmodern,amsmath,amssymb,amsthm,mathtools,booktabs,microtype}
\usepackage[margin=27mm]{geometry}
\usepackage[hidelinks]{hyperref}
\newtheorem{theorem}{Theorem}[section]
\newtheorem{lemma}[theorem]{Lemma}
\newtheorem{proposition}[theorem]{Proposition}
\newtheorem{corollary}[theorem]{Corollary}
\newtheorem{remark}[theorem]{Remark}
\newcommand{\Z}{\mathbb Z}
\newcommand{\R}{\mathbb R}
\newcommand{\C}{\mathbb C}
\newcommand{\norm}[1]{\lVert#1\rVert}
\newcommand{\im}{\operatorname{im}}
\newcommand{\Rea}{\operatorname{Re}}
\newcommand{\e}{\mathrm e}
\providecommand{\SourceCommit}{\errmessage{Supply the frozen SourceCommit}}
\pdfinfoomitdate=1\pdftrailerid{}
\ifdefined\pdfsuppressptexinfo\pdfsuppressptexinfo=15\fi
\title{Coverage and coexactness in checkerboard slice observations}
\author{Lluis Eriksson\\\small THE ERIKSSON PROGRAMME}
\date{12 November 2026}
\begin{document}
\maketitle
\begin{abstract}
We study original-channel rectangular means for a finite three-dimensional
checkerboard gradient energy with weak coefficient $c/48$ and strong coefficient
$cn^2$. Each partition has at most $36n$ rows and is specified by one cut per
slice. If the cut diagonal is a permutation, actual singleton observations give
an explicit logarithmic coercivity bound and a finite quadratic penalty, valid
on all fields and therefore on the actual curl-adjoint image. For arbitrary
cuts, the full means exclude every nonzero field of exactly zero strong energy;
this qualitative fact does not restore the target coefficient $c$. Conversely,
at an explicit, extremely large dyadic threshold, every cut diagonal with image
of size at most $n/32$ admits a real, nonzero, genuinely coexact invisible field
of energy below $c/2$ times its squared norm. We also prove a contrast-uniform
factor-four energy bound for the true coexact projection, give an explicit
counterexample to its preservation of these mean kernels, and quantify a
necessary correction cost. The last estimate concerns the correction alone,
not the energy of the corrected field. The permutation case at the original
target for the nonlinear example $F(i)=i+2i^2$, historical priority, and a
matching logarithmic rate remain open; some affine permutation subfamilies
already have counterexamples at that target.
These are finite analytical results, not formal verification or a physical
Yang--Mills spectral-gap theorem.
\end{abstract}

\section{Problem, scope and principal conclusions}
Low-dimensional observations of a high-contrast energy depend both on the
observed functionals and on the domain on which a restoring inequality is
required. Here the domain distinction is concrete: a field in the actual
curl-adjoint image need not remain there after rectangular averaging, while
projection into that image need not preserve rectangular mean zeros. We keep
these two operations separate throughout.

The purpose is a consolidated account of the programme's finite analytical
sources 266--273. Their complete proofs and source-specific reports accompany
the reproducibility package. The preceding checkerboard working paper
\cite{previous} treats broader observation-count and tagged-window questions.
We do not claim a new general theory of capacity or Hodge decomposition.
Logarithmic point estimates in two dimensions are classical
\cite{lawler}; stability of projections in differential complexes is also an
established subject \cite{afw}. Every estimate needed for this particular
coefficient family is derived below. A bounded comparison with these primary
sources does not establish historical originality or wider practical relevance.

The results describe three distinct situations. Permutation diagonals yield a
paid lower bound of order $c/(1+\log n)$. Arbitrary diagonals still exclude
exactly strong-null invisible fields. Low-coverage diagonals nevertheless fail
at target $c$ through fields of positive, small strong energy. Thus qualitative
injection on the strong kernel is insufficient, and coverage alone is not
proved sufficient. We establish no condition-number or solver guarantee.

\subsection{Finite model and actual observations}
Let $n\geq4$ be even, $X=(\Z/n\Z)^3$, $V=n^3$, and $c>0$. Set
\[
 a=c/48,\qquad H=cn^2,\qquad S(x)=x_1+x_2+x_3\pmod n.
\]
The two original channels are orthogonal. Each has three spatial components.
For $\partial_j f(x)=f(x+e_j)-f(x)$, channel 1 has coefficient $H$ at odd
source parity and $a$ at even source parity; channel 2 has the reverse
assignment. The coefficient is scalar across spatial directions and
components within a fixed source and channel. Write
\[
 Q(u)=Q_a(u)+Q_H(u)
 =\sum_{\ell=1}^2\sum_{x,j,r} b_\ell(x)
       |\partial_j u_{\ell,r}(x)|^2.
\]
All norms are counting norms. All forward labels are retained.

For a one-form $v$, the forward curl is
$(d_1v)_{jl}=\partial_jv_l-\partial_lv_j$, $j<l$. The domain of interest is
\[
 U=(\im d_1^*)^2.
\]
This is the actual image with counting adjoints; constants are excluded.
We never replace it by a divergence-free space without removing harmonic modes.

Choose any $t_i\in\{0,\ldots,n-1\}$. On the slice $x_1=i$, partition $x_2$
into the nonempty intervals
\[
 [0,t_i-1],\quad\{t_i\},\quad[t_i+1,n-1],
\]
and partition $x_3$ into $\{0\}$ and $[1,n-1]$. Every original spatial
component in each channel is observed on every resulting cell. A row is
$V^{-1/2}|B|^{-1/2}\sum_{x\in B}u_{\ell,r}(x)$.
Call the full map $C_t$, and put $F(i)=i+t_i\pmod n$. If $Z_0$ and $Z_*$
count cuts equal to $0$ and $n-1$, respectively, then the raw row count is
\[
 36n-12(Z_0+Z_*),\qquad 24n\leq\#\text{rows}\leq36n.
\]
Every cell is a nonwrapping chart rectangle. No cyclic-union mean or mixed
channel replaces these rows. In particular $(i,t_i,0)$ is always an actual
singleton. The target is
\begin{equation}\label{target}
 Q(u)/V+\Phi(C_tu)\geq c\norm{u}^2/V,\qquad \Phi(0)=0,
\end{equation}
on $U$, with a finite observation penalty. Our counterexamples also defeat
quadratic penalties. We do not assert failure for every canonical partition.

\section{The strong graph and a paid point estimate}\label{green}
Fix a scalar component in a channel with strong-source parity $p$. Its strong
components are
\[
 G_s=\{x:S(x)\in\{s,s+1\}\},\qquad s\equiv p\pmod2.
\]
Indeed the undirected strong path
$x\to x+e_j\to x+e_j-e_k$ connects source sites by zero-sum generators.
The vectors $e_1-e_3,e_2-e_3$ generate every zero-sum displacement modulo $n$;
each target has a rung. No strong edge leaves a paired layer. Each component
therefore has $2n^2$ vertices.

Parametrize its source by $x_s(y)=(y_1,y_2,s-y_1-y_2)$ and its target by
$x_s(y)+e_3$, $y\in(\Z/n\Z)^2$. For values $b(y),d(y)$ the exact component
energy is
\[
 E_s=\sum_y|b(y)-d(y)|^2+
              \sum_{y,j=1,2}|b(y)-d(y+e_j)|^2.
\]
Thus $Q_H=H\sum E_s$, summing channels and spatial components. In the
unitary two-dimensional Fourier transform, with normalization $n^{-1}$,
the Hermitian matrix has diagonal 3 and off-diagonals $-A,-\overline A$,
where $A=1+\e^{i\theta_1}+\e^{i\theta_2}$. Its eigenvalues are $3\pm|A|$,
and
\begin{equation}\label{det}
 9-|A|^2=|1-\e^{i\theta_1}|^2+|1-\e^{i\theta_2}|^2+
                                      |\e^{i\theta_1}-\e^{i\theta_2}|^2.
\end{equation}
At zero frequency the eigenvalues are $0,6$; only the constant on both
layers is omitted below.

\begin{lemma}\label{point}
Put $L=1+\log n$. On one actual strong component, for any $f$ of total mean
zero and any vertex $v$,
\[
 \norm{f}^2\leq(3n^2/8)E_s(f),\qquad
 |f(v)|^2\leq G_nE_s(f),\qquad G_n\leq2L,
\]
where the actual inverse diagonal is
\begin{equation}\label{G}
 G_n=\frac1{n^2}\left[\frac1{12}+
              \sum_{k\ne0}\frac3{9-|A(k)|^2}\right].
\end{equation}
\end{lemma}
\begin{proof}
Choose centered residues $r_j\in[-n/2,n/2-1]$. The finite sine bound gives
$|1-\e^{2\pi i r_j/n}|^2\geq16r_j^2/n^2$. Equation~\eqref{det} and
$3+|A|\leq6$ imply $3-|A|\geq8/(3n^2)$ at nonzero frequency. The
zero-frequency antisymmetric mode also obeys this gap, proving the norm bound.

Diagonalize the finite positive-energy coordinates and apply weighted
Cauchy--Schwarz to point evaluation. At nonzero frequency the inverse matrix
diagonal is $3/(9-|A|^2)$. The remaining zero-frequency mode contributes
$(1/2)/6=1/12$. Squared Fourier normalization is $1/n^2$, proving~\eqref{G}.
Discarding only the third nonnegative term in~\eqref{det} gives
\[
 G_n\leq\frac1{12n^2}+\frac3{16}
                \sum_{r\ne0}(r_1^2+r_2^2)^{-1}.
\]
The shell $\max(|r_1|,|r_2|)=j$ has at most $8j$ points and denominator
at least $j^2$. Consequently the sum is at most
$8\sum_{j=1}^{n/2}j^{-1}\leq8(1+\log(n/2))$, and $G_n\leq2L$ follows.
The possibly smaller Nyquist shell only improves this finite upper bound.
\end{proof}

\section{Permutation coverage gives logarithmic coercivity}
\begin{theorem}\label{lower}
If $F$ is a permutation, then for every ambient field $u$,
\begin{equation}\label{penalty}
 Q(u)/V+\frac{4cn^2}{9L}\norm{C_tu}^2
                  \geq\frac{c}{9L}\norm{u}^2/V.
\end{equation}
On the actual observation kernel there is the stronger estimate
\begin{equation}\label{kernel}
 C_tu=0\quad\Longrightarrow\quad
                  Q(u)\geq\frac{c}{5L}\norm{u}^2.
\end{equation}
Both hold on $U$ by restriction, without invariance of rectangular averaging.
\end{theorem}
\begin{proof}
For each channel and strong component select the unique actual singleton
$v_s=(i_s,t_{i_s},0)$ for which $F(i_s)=s$ is its source residue. Let a
scalar component restriction have mean $m$ and residual $r=f-m$. If the
observations vanish, $m=-r(v_s)$. Lemma~\ref{point} and orthogonal counting
decomposition yield
\[
 \norm{f}^2=\norm r^2+2n^2|m|^2
 \leq[3/8+4L]n^2E_s(f)\leq5n^2LE_s(f).
\]
Sum the original scalar components and use $Q\geq Q_H=cn^2\sum E_s$.

For an arbitrary field, $m=f(v_s)-r(v_s)$, so
\[
 \norm f^2\leq4n^2|f(v_s)|^2+[3/8+8L]n^2E_s(f)
           \leq4n^2|f(v_s)|^2+9n^2LE_s(f).
\]
The selected original singleton rows are distinct, and their squared
unnormalized values sum to at most $V\norm{C_tu}^2$. Therefore
$\norm u^2\leq4n^2V\norm{C_tu}^2+(9L/c)Q(u)$, which is~\eqref{penalty}.
This pays the finite-penalty step directly on all fields.
\end{proof}

For dyadic $n$, the nonlinear example $t_i=2i^2\pmod n$ has permutation
diagonal: $F(i)-F(j)=(i-j)[1+2(i+j)]$, and the bracket is odd and invertible.
The map is nonaffine for $n\geq8$; this elementary example is not claimed
new. Theorem~\ref{lower} does not attain target $c$, prove logarithmic
necessity, or decide~\eqref{target} for this example.

\section{Full means exclude exactly strong-null invisible fields}
\begin{lemma}\label{midpoint}
If every adjacent midpoint of a finite real cyclic sequence
$A_0,\ldots,A_{M-1}$ belongs to its range, the sequence is constant.
\end{lemma}
\begin{proof}
The case $M=1$ is immediate. Otherwise choose $\sigma(j)$ with
$2A_{\sigma(j)}-A_j-A_{j+1}=0$, cyclically. The integer coefficient matrix
has row sums zero. Modulo 2 it is $I+T$, where $T$ is the next-coordinate
shift. Keep rows $0,\ldots,M-2$ and columns $1,\ldots,M-1$: the resulting
matrix modulo 2 is lower bidiagonal with diagonal entries 1. Its integer
minor is odd and nonzero, so real rank is at least $M-1$. Constants already
belong to its kernel, so that kernel consists only of constants. This does
not require rational values or an injective selection $\sigma$.
\end{proof}

\begin{theorem}\label{exactzero}
For every cut map $t$ in the stated slice family,
$\ker C_t\cap\ker Q_H=\{0\}$ on the whole real ambient space, hence on $U$.
No parity, permutation or motion assumption on the cuts is required.
\end{theorem}
\begin{proof}
Fix a scalar component in a channel of strong-source parity $p$.
Zero strong energy makes it $f(S)$ with $f(s)=f(s+1)$ for $s\equiv p$,
by the actual graph connectivity in Section~\ref{green}. All cell sums zero
give global sum zero and hence $\sum_s f(s)=0$. Choose a periodic primitive
$G(s+1)-G(s)=f(s)$. Strong locking implies
$G(s+1)=[G(s)+G(s+2)]/2$ when $s\equiv p$.

The actual singleton gives $f(F(i))=0$. The actual left cell, or the
identity when $t_i=0$, gives
\[
 0=\sum_{j=0}^{t_i-1}f(i+j)=G(F(i))-G(i).
\]
This remains valid at $t_i=n-1$; no wrapping cell is introduced. For any
$i\equiv1-p$, if $F(i)\equiv p$, then $G(i)$ is already a value on parity
$p$. Otherwise $f(F(i))=0$ makes $G(F(i))=G(F(i)+1)$, also a value on
parity $p$. Thus every value of $G$ on parity $1-p$ belongs to its range
on parity $p$. These are precisely the adjacent midpoints of
$A_j=G(p+2j)$, $0\leq j<n/2$. Lemma~\ref{midpoint} makes $G$, and then
$f$, constant and zero, respectively. Apply this to every original component.
\end{proof}

This theorem excludes one exact-zero mechanism, not a small positive energy.
It supplies no quantitative size-uniform bound and does not remove the
permutation assumption from Theorem~\ref{lower}.

\section{Low coverage admits a true coexact invisible field}\label{lowsection}
\begin{theorem}\label{low}
Let $n=2^m\geq1024$, $K=|\im F|$, and suppose
\[
 K\leq n/32,\qquad T=\log(n/64)\geq2^{34}.
\]
There is a real nonzero $z\in U$ with $C_tz=0$ and
$Q(z)<(c/2)\norm z^2$. Consequently no finite penalty vanishing at zero
restores target $c$ for this map. Within this slice family and at these sizes,
target $c$ requires $K>n/32$; neither sufficiency nor optimality is asserted.
\end{theorem}

We give the full construction. Write $\delta_j=I-\tau_{-e_j}$, and let
$\delta=I-\tau_{-1}$ on a profile of $S$. For a scalar potential $\Psi$ set
\[
 \mathcal P\Psi=(0,\delta_2\delta_3^2\Psi,-\delta_2^2\delta_3\Psi).
\]
This has the actual two-form preimage
$A_{23}=\delta_2\delta_3\Psi$, $A_{12}=A_{13}=0$, because
$(d_1^*A)_2=\delta_3A_{23}$ and $(d_1^*A)_3=-\delta_2A_{23}$.
Use this field in channel 1 and zero in channel 2. Every original rectangular
sum telescopes to four endpoint values of $\delta_3\Psi$ or
$-\delta_2\Psi$. It suffices to annihilate both gradients at
\begin{equation}\label{tags}
 x_1=i,\qquad x_2\in\{-1,t_i-1,t_i\},\qquad x_3\in\{-1,0\}.
\end{equation}
Duplicated tags at endpoint cuts only remove requirements.

\subsection{A derived profile kernel with a uniform inverse}
Put $B=\im F$, $Z=B+\{-3,-2,-1,0\}$, so $|Z|\leq4K$.
Use unitary one-dimensional characters $e_p(s)=n^{-1/2}\e^{2\pi ips/n}$.
For $r=0,\ldots,4K$, let $p_r=n/4+r$, $h=n/2$,
$\zeta_r=\e^{2\pi ip_r/n}$ and $b(\zeta)=1-\zeta^{-1}$. Define
\begin{align*}
 \phi_r&=(1+\zeta_r)e_{p_r}+(1-\zeta_r)e_{p_r+h},\\
 \psi_r&=\frac{1+\zeta_r}{b(\zeta_r)^3}e_{p_r}
           +\frac{1-\zeta_r}{b(-\zeta_r)^3}e_{p_r+h}.
\end{align*}
The $4K+1$ first frequencies lie in $[n/4,3n/8]$, including the endpoint
when $K=n/32$; partners lie in $[3n/4,7n/8]$. They are distinct and nonzero.
The phases lie in $[\pi/2,3\pi/4]$, so both multipliers have modulus at
least $1/2$. Disjoint supports give independence and, for every combination,
\begin{equation}\label{inverse}
 \phi=\delta^3\psi,\qquad \norm\psi_S^2\leq64\norm\phi_S^2.
\end{equation}
Evaluation on $Z$ has at most $4K$ outputs. Rank--nullity therefore supplies
a nonzero $\psi$ in this $4K+1$ dimensional space with $\psi|_Z=0$.
The nonzero multipliers make $\phi\ne0$. No conditioning or coefficient bound
for the evaluation matrix is assumed.

For odd $s$, each $\phi_r$ has zero forward difference, since its parity-pair
bracket is
$(\zeta-1)(1+\zeta)-(-\zeta-1)(1-\zeta)=0$. Thus
$z_0=(0,\phi(S),-\phi(S))=\mathcal P[\psi(S)]$ is exactly strong locked,
and
\begin{equation}\label{base}
 \norm{z_0}^2=2n^2\norm\phi_S^2>0.
\end{equation}

\subsection{Finite cutoff and every original mean}
For $y=(x_2,x_3)$ put $r(y)=\max(d_n(x_2,0),d_n(x_3,0))$, with cyclic
integer distance, and $R=n/8$. Define
\[
 \chi(y)=\begin{cases}0&r\leq8,\\ \log(r/8)/T&8<r<R,\\1&r\geq R.
 \end{cases}
\]
Let $E_\chi=\sum_{y,j=2,3}|\chi(y+e_j)-\chi(y)|^2$. A nonzero edge
crosses radii $j,j+1$, $8\leq j<R$, with difference at most $1/(Tj)$.
The shell has $8j$ points. Counting the two outgoing labels from both shells
bounds relevant edges by $16j+16(j+1)\leq48j$. Hence
\begin{equation}\label{cap}
 E_\chi\leq\frac{48}{T^2}\sum_{j=8}^{R-1}\frac1j\leq\frac{96}{T},
 \quad \sum_y\chi^2\geq\frac{3n^2}{4},
 \quad\norm{\delta^\beta\chi}_{2D}^2\leq4^{|\beta|-1}E_\chi
 \quad(|\beta|\geq1).
\end{equation}
The last bound retains one first difference and bounds each other difference
by operator norm 2. Backward and forward first norms agree by reindexing.
The mass estimate discards the square of radius $R$, whose size
$(2R+1)^2\leq n^2/4$. Torus seam edges are outside the varying region.

Take $\Psi=\psi(S)\chi$ and $z=\mathcal P\Psi$. At fixed tags
$x_2=-1,x_3=-1,0$, all potential values used by the two first gradients have
transverse radius at most 2, so vanish. At moving tags
$x_2=t_i+u,x_3=v$, $u,v\in\{-1,0\}$, the values and their backward
neighbors have sum coordinates in $F(i)+\{-3,-2,-1,0\}\subset Z$.
They vanish irrespective of the cutoff. Thus~\eqref{tags} holds and every
original normalized cell mean is zero. The displayed preimage already pays
actual coexactness; no divergence converse is used.

\subsection{Nonzero norm and original source energy}
Write $z=\chi z_0+E$. For either order-three derivative, the exact backward
Leibniz formula has terms
\[
 \binom\alpha\beta\bigl[\tau_{-\beta}
       \delta^{\alpha-\beta}\psi(S)\bigr]\delta^\beta\chi.
\]
Beyond $\beta=0$ their total multiplicity is 7 per nonzero component.
At fixed transverse $y$, the $x_1$ sum of any shifted profile square is its
one-dimensional norm. A term with $b=|\beta|\geq1$ has squared norm at most
$4^{3-b}\norm\psi_S^2\,4^{b-1}E_\chi=16\norm\psi_S^2E_\chi$.
Seven-term Cauchy and two components, then~\eqref{inverse}--\eqref{base}, give
\begin{equation}\label{err}
 \norm E^2\leq2^{11}\norm\psi_S^2E_\chi
       \leq\frac{2^{16}E_\chi}{n^2}\norm{z_0}^2.
\end{equation}
This uses the norm of the entire profile, not a sum of absolute Fourier
coefficients; there is no factor depending on $K$. The mass bound gives
$\norm{\chi z_0}\geq3\norm{z_0}/4$. Our threshold implies
$E_\chi<2^{-26}<1$; $n\geq1024$ then gives $\norm E\leq\norm{z_0}/4$ and
\begin{equation}\label{nz}
 \norm z^2\geq\norm{z_0}^2/4>0.
\end{equation}

At odd sources the strong difference of $\chi z_0$ is just a cutoff
difference times $z_0(x)$. Summing profiles in $x_1$ and dropping only the
odd restriction yields
$Q_H(\chi z_0)\leq2H\norm\phi_S^2E_\chi=cE_\chi\norm{z_0}^2$.
The full unweighted three-direction gradient is at most 12 times counting
norm, so $Q_H(E)\leq12H\norm E^2<2^{20}cE_\chi\norm{z_0}^2$.
The ambient positive quadratic form and~\eqref{nz} give
\[
 Q_H(z)<2^{22}cE_\chi\norm{z_0}^2
         \leq2^{24}cE_\chi\norm z^2<(c/4)\norm z^2.
\]
Also $Q_a(z)\leq12a\norm z^2=(c/4)\norm z^2$. Thus the complex field
has quotient below $c/2$. Real and imaginary potential parts are actual
coexact invisible fields; their energies and norm squares add. Some nonzero
real part has no larger quotient. This proves Theorem~\ref{low}.
Theorem~\ref{exactzero} implies its strong energy is strictly positive;
it was not used to construct the witness or its upper bound.

The threshold is enormous and conservative, with no practical-size claim.
For example $t_i=((i\bmod2)-i)\bmod n$ has $F(i)=i\bmod2$, is nonaffine,
and has exactly $36n-24$ full rows. Only $i=0,1$ have zero cuts, and none
has cut $n-1$. It is covered by Theorem~\ref{low}. Permutations, including
$i+2i^2$, are outside its low-coverage hypothesis.

\section{Energy stability of the actual coexact projection}\label{projection}
Let $P$ be counting orthogonal projection onto $U$, independently in the two
original channels. At nonzero physical Fourier frequency $k$ put
$g_j=\e^{2\pi i k_j/n}-1$. Its symbol is
\begin{equation}\label{P}
 P_g=I-gg^*/|g|^2,\qquad P_0=0.
\end{equation}
Indeed the curl symbol is $g\wedge$, and for $g^*v=0$ the actual two-form
$A_{jl}=(g_jv_l-g_lv_j)/|g|^2$ has adjoint $v$. Conversely antisymmetry
gives $g^*d_1^*=0$. At zero frequency the image is zero. Symbols conjugate
at opposite frequencies, so $P$ preserves real fields.

\begin{theorem}\label{stableP}
For every ambient field, $Q_H(Pu)\leq4Q_H(u)$ and
$Q_a(Pu)\leq4Q_a(u)$. The same conclusion holds for any nonnegative scalar
parity weights $a,H$, with constant independent of contrast and $n$.
\end{theorem}
\begin{proof}
For nonzero $g,h\in\C^3$ and spatial vectors $A,B$, write
$R_j=g_jA-h_jB$ and $s^2=1-|g^*h|^2/(|g|^2|h|^2)$. The difference of
the two rank-one complementary projections has norm $s$. To verify this
over $\C$, adjust the phase of $h/|h|$ so its inner product with $g/|g|$
is $\alpha\geq0$, and write it $\alpha e+sv$. On their two-dimensional
span the rank-one projection difference is
$\left(\begin{smallmatrix}s^2&-\alpha s\\-\alpha s&-s^2\end{smallmatrix}\right)$,
whose square is $s^2I$; it is zero on the orthogonal complement. Parallel
vectors give zero difference directly.

Completing the square, with the Hermitian scalar $g^*h$, gives
\[
 \sum_j\norm{R_j}^2=|g|^2\norm{A-(g^*h)B/|g|^2}^2
                           +|h|^2s^2\norm B^2.
\]
The projected residual is
$R'_j=P_gR_j+h_j(P_g-P_h)B$. Each tuple on the right has norm at most
$\norm R$, so $\norm{R'}\leq2\norm R$. Replacing $B$ by $-B$ gives the
same estimate for the plus residual. If either symbol is zero, the actual
$P_0=0$ reduces the estimate to norm contraction of the other projection.

Pair physical frequencies $k,k+\eta$, $\eta=(n/2,n/2,n/2)$, once each.
Their forward symbols are $g,h$, with $h_j=-\e^{2\pi i k_j/n}-1$.
On an odd source their gradient tuple is $e_k(x)(g_jA-h_jB)$; on even
sources it is $e_k(x)(g_jA+h_jB)$. Distinct unordered pairs are orthogonal
on each parity sublattice, since its character sum vanishes unless the
frequency difference is $0$ or $\eta$. The restricted squared character
norm is exactly $1/2$. Each parity energy is therefore half the sum of
these squared tuples. Apply the residual estimate and the nonnegative
original weight. The reversed channel only exchanges the two signs.
\end{proof}

Norm contraction supplies no lower bound for $\norm{Pu}$. The theorem also
does not preserve any observation row or assert energy commutation. The next
example makes the remaining mean obstruction explicit.

\section{Loss of means and a necessary correction cost}
\begin{proposition}\label{loss}
For $4\mid n$, the real ambient field in channel 1
\[
 u(x)=-2\sin(\pi x_1/2)\sin(\pi x_3/2)e_1
\]
has $C_tu=0$ for every cut map, but
$(Pu)_3(0,t_0,0)=-1/2$. Thus its true coexact projection is visible for
every map in the specified full slice family.
\end{proposition}
\begin{proof}
The field is independent of $x_2$, vanishes at $x_3=0$, and its sine sum
over $[1,n-1]$ is zero. Every original cell mean therefore vanishes.
Use the complex precursor
$w=\e^{i\pi x_1/2}(\e^{i\pi x_3/2}-\e^{-i\pi x_3/2})e_1$, whose real
part is $u$. At its two frequencies the symbols are
$g_+=(i-1,0,i-1)$ and $g_-=(i-1,0,-i-1)$, both squared norm 4.
At $x_3=0$, the first two components of
$P_{g_+}e_1-P_{g_-}e_1$ vanish, and the third is
\[
 -[(i-1)-(-i-1)]\overline{(i-1)}/4=(-1+i)/2.
\]
The operator $P$ is real, so at $x_1=0$ its real third value is $-1/2$,
independent of $x_2$. The actual normalized singleton row is
$-1/(2\sqrt V)\ne0$, including endpoint cuts.
\end{proof}

\begin{theorem}\label{repair}
For any permutation diagonal and any $w\in U$, with channel-1 components
$w_1,w_2,w_3$, select $v_s$ as in Theorem~\ref{lower} for odd $s$. Then
\begin{equation}\label{longtrace}
 H\sum_{s\ \mathrm{odd}}|w_1(v_s)+w_2(v_s)+w_3(v_s)|^2
                                                \leq6LQ_H(w).
\end{equation}
For the input in Proposition~\ref{loss}, every coexact $r$ satisfying
$C_t(Pu+r)=0$ obeys
\begin{equation}\label{repaircost}
 Q_H(r)\geq\frac{c\norm u^2}{48L},\qquad \norm u^2=n^3.
\end{equation}
Equation~\eqref{repaircost} concerns $r$ only, not the final field $Pu+r$.
\end{theorem}
\begin{proof}
For a true coexact field, $d_0^*w=0$ follows from $d_1d_0=0$, and every
spatial component has global sum zero by its actual backward-difference
adjoint formula. Let $W_j(s)=\sum_{S(x)=s}w_j(x)$ and
$T(s)=\sum_jW_j(s)$. Summing divergence on a layer yields
$T(s)-T(s-1)=0$. Its global sum is zero, hence $T(s)=0$.
Thus $f=w_1+w_2+w_3$ has zero mean on every individual sum layer and on
every strong pair. No converse divergence characterization is needed.

Lemma~\ref{point} gives $|f(v_s)|^2\leq G_nE_s(f)$, and edgewise
Cauchy gives $E_s(f)\leq3\sum_jE_s(w_j)$. The selected points are in
distinct strong components. Summation and $G_n\leq2L$ give~\eqref{longtrace}.
The other original channel contributes nonnegative energy and is not mixed
into the trace. The spatial sum is a derived combination of three actual
rows, with its factor 3 paid.

At every $x_1=i,x_3=0$, the first two components of $Pu$ vanish and its
third component is
$\Rea[\e^{i\pi i/2}(-1+i)/2]\in\{-1/2,1/2\}$.
Any repair must have $r(v_s)=-Pu(v_s)$, separately in each original spatial
row. There are $n/2$ selected points, so the squared spatial-sum trace of
$r$ totals $n/8$. Equation~\eqref{longtrace} gives
$Q_H(r)\geq Hn/(48L)=cn^3/(48L)$. Finally
$\norm u^2=4n(\sum_i\sin^2(\pi i/2))^2=n^3$.
\end{proof}

This excludes a strong-null coexact correction for this input. It excludes
neither paid corrections nor cancellation of their energy against $Pu$.
For example $r=-Pu$ is admissible and leaves final field zero. We supply
no nonzero corrected invisible witness here.

\section{Remaining questions and review context}
The original coefficient $c$ for the nonlinear permutation $F(i)=i+2i^2$
remains undecided. Theorem~\ref{lower} is a logarithmic lower bound;
Theorem~\ref{low} is a counterexample for a different coverage regime.
Neither establishes necessity of a logarithmic observation count or a
matching logarithmic kernel rate for permutation diagonals. The projection
and repair results explain why an ambient construction cannot simply be
transferred into the actual domain, but they do not close that gap.

All results concern finite counting models. There is no continuum limit,
physical spectral-gap assertion, practical threshold, optimal coefficient,
solver guarantee or claim that the finite estimates solve Yang--Mills.
The programme's separate PR29 record remains 20/41 with TermSource $=0$.

\subsection*{Model assessments and review context}
The mathematical source notes received complete, exact-source analytical
checks by an involved Codex/OpenAI assistant with shared author/programme and
prior-audit context. The exact served backend, version and effort were not
exposed and are not inferred. These checks are not blind, independently
established peer review, formal proofs or numerical experiments. Their full
reports, assignments, objections and read scopes are retained in the source
package; absence of a detected defect is not an originality finding.

An exact-PDF assessment, if supplied in the deposit metadata, is a separate
artifact tied to that PDF's hash. No numerical assessment of an earlier
working paper is transferred to this manuscript. Source checks, PDF/model
assessment, deposit, working-paper publication, human peer review and editorial
acceptance are distinct statuses. The author and programme control the
manuscript. The author is also the owner, creator and administrator of AIRR,
the intended working-paper destination; this author--operator role is disclosed
and does not constitute independent editorial or human peer review.

\subsection*{Source and reproducibility}
The production source commit is
\texttt{\SourceCommit}. The accompanying package retains original sources
266--273 and their full reports with raw-byte SHA-256 manifests. This manuscript
contains the finite definitions and proofs needed for its statements. Linux
production logs and exact-PDF/render hashes, when generated, are separate
evidence; a successful existing-core Lean/oracle CI run does not formally verify
these plain-text analytical proofs. No source-stage claim of numerical or
formal verification is made.

\begin{thebibliography}{9}
\bibitem{lawler} G. F. Lawler and V. Limic,
\emph{Random Walk: A Modern Introduction}, Cambridge University Press, 2010.
Author-hosted draft: \url{https://www.math.uchicago.edu/~lawler/srwbook.pdf}.
Classical background, particularly two-dimensional potential theory; our finite
Green normalization is derived in Section~\ref{green}.
\bibitem{afw} D. N. Arnold, R. S. Falk and R. Winther,
Finite element exterior calculus: from Hodge theory to numerical stability,
\emph{Bulletin of the American Mathematical Society} 47 (2010), 281--354.
\url{https://arxiv.org/abs/0906.4325}.
The cochain projections there are not identified with the projection in
Section~\ref{projection}.
\bibitem{previous} L. Eriksson,
\emph{Canonical and unrestricted observations in a finite checkerboard gradient
model}, AIRR working paper ARR-2026-597T88D7WZ98HR0Z, version 2, 2026.
\url{https://airr.science/papers/ARR-2026-597T88D7WZ98HR0Z/versions/v2/}.
Working-paper status and its own assessments do not imply acceptance or carry
over to the present PDF.
\end{thebibliography}
\end{document}
